\documentclass[11pt,letterpaper]{article}
\usepackage[margin=1in]{geometry}
\usepackage[utf8]{inputenc}
\usepackage{parskip}
\usepackage{hyperref}
\usepackage{enumitem}

\hypersetup{
    colorlinks=true,
    linkcolor=black,
    urlcolor=blue,
    pdftitle={Sport Field and Sport Court Strategy Submission},
    pdfauthor={Michael Kernaghan}
}

\begin{document}

\begin{center}
{\Large \textbf{Submission: Include Croquet in Sport Field and Sport Court Strategy}}
\end{center}

\noindent
\textbf{To:} Cory Nesci, Planning Manager \\
\textbf{Email:} fieldsandcourts@abbotsford.ca \\
\textbf{Re:} Sport Field and Sport Court Strategy -- Public Feedback \\
\textbf{Date:} February 26, 2026

\vspace{0.5cm}
\noindent\rule{\textwidth}{0.4pt}
\vspace{0.5cm}

\section*{Executive Summary}

I am writing to request that \textbf{croquet} be explicitly included in the City of Abbotsford's Sport Field and Sport Court Strategy. While lawn bowling, bocce, and horseshoe pitches are currently recognized under ``Other outdoor sport infrastructure,'' croquet---a similar lawn sport with an active local presence---is notably absent from the strategy.

\textbf{Importantly, the Abbotsford Croquet Club currently shares facilities with the Abbotsford Lawn Bowling Club.} This means croquet is already part of the city's sport infrastructure, yet remains unrecognized in the long-term planning strategy.

\section*{Background: Croquet in Abbotsford}

The Abbotsford Croquet Club is an established organization that has been serving the community for years. \textbf{The club currently shares facilities with the Abbotsford Lawn Bowling Club}, demonstrating successful multi-use of city sport infrastructure.

Croquet is:

\begin{itemize}[leftmargin=*]
    \item \textbf{An organized sport} with regular club play, tournaments, and competitive opportunities
    \item \textbf{Already using city facilities}---the shared lawn bowling/croquet facility represents existing city investment
    \item \textbf{Growing in the region}---part of the Pacific Northwest croquet circuit including clubs in Victoria, Vancouver, Seattle, and Portland
    \item \textbf{Accessible to diverse populations}---appeals to adults of all ages, requires minimal equipment, and has low barriers to entry
    \item \textbf{Social and strategic}---combines physical activity with tactical thinking and community building
\end{itemize}

\section*{Why Croquet Deserves Recognition}

\subsection*{1. Consistency with Current Strategy}
The strategy already includes lawn bowling under ``Other outdoor sport infrastructure.'' \textbf{The Abbotsford Croquet Club currently shares the same facilities with the Lawn Bowling Club}, yet only lawn bowling is recognized in the strategy. This omission is inconsistent because:

\begin{itemize}[leftmargin=*]
    \item \textbf{Same facilities}---croquet uses the same city-maintained grass surfaces as lawn bowling
    \item \textbf{Same infrastructure}---the city is already supporting croquet through its lawn bowling facility investment
    \item \textbf{Comparable demographics}---both sports appeal to older adults and retirees
    \item \textbf{Similar community benefits}---social connection, outdoor recreation, accessible physical activity
    \item \textbf{Proven compatibility}---the existing shared-use arrangement demonstrates successful multi-sport facility management
\end{itemize}

\subsection*{2. Alignment with City Goals}
Including croquet supports the strategy's stated objectives:
\begin{itemize}[leftmargin=*]
    \item \textbf{Accommodating population growth}---particularly the aging demographic expected over the next 15 years
    \item \textbf{Improving facility management}---croquet can share spaces with other lawn sports
    \item \textbf{Collaborative opportunities}---potential partnerships with the Abbotsford Croquet Club for programming and maintenance
\end{itemize}

\subsection*{3. Infrastructure Already in Place}
\textbf{The city is already supporting croquet through the shared lawn bowling facility.} This means:
\begin{itemize}[leftmargin=*]
    \item \textbf{Zero additional infrastructure cost}---the facility already exists and serves both sports
    \item \textbf{Proven multi-use model}---the current arrangement demonstrates effective facility sharing
    \item \textbf{Efficient use of city investment}---recognizing croquet acknowledges the full value of existing infrastructure
    \item \textbf{Scalability}---the successful shared-use model could be replicated in future facility planning
\end{itemize}

\subsection*{4. Regional and Competitive Context}
The Pacific Northwest has an active croquet network. Abbotsford's recognition of croquet in its long-term strategy would:
\begin{itemize}[leftmargin=*]
    \item Position the city as a regional destination for tournaments
    \item Attract visitors for multi-day events
    \item Support local players competing at provincial and national levels
\end{itemize}

\section*{Specific Requests}

I respectfully request that the City:

\begin{enumerate}[leftmargin=*]
    \item \textbf{Add croquet} to the list of sports explicitly recognized under ``Other outdoor sport infrastructure''
    \item \textbf{Include croquet facilities} in the inventory assessment of current and potential future sites
    \item \textbf{Consult with the Abbotsford Croquet Club} during Stage 4 finalization to understand facility needs and opportunities
    \item \textbf{Consider multi-use lawn sport areas} that can accommodate croquet, lawn bowling, and bocce to maximize facility utilization
\end{enumerate}

\section*{Conclusion}

The omission of croquet from the Sport Field and Sport Court Strategy appears to be an oversight rather than a deliberate exclusion. \textbf{Croquet is already part of Abbotsford's sport infrastructure through the shared lawn bowling facility}---the strategy simply fails to recognize it.

As the City plans for the next 15+ years of sport infrastructure, formally recognizing croquet alongside lawn bowling will:
\begin{itemize}[leftmargin=*]
    \item \textbf{Accurately reflect current facility use}---the strategy should document what's already happening
    \item \textbf{Ensure comprehensive planning} for diverse recreational needs
    \item \textbf{Support an existing and growing local club} that's already using city facilities
    \item \textbf{Validate the multi-use model} that's proven successful at the lawn bowling/croquet facility
    \item \textbf{Maximize the value of grass field investments} by acknowledging all user groups
\end{itemize}

I am happy to provide additional information about croquet's local presence, facility requirements, or community benefits. Thank you for considering this feedback as you finalize the Sport Field and Sport Court Strategy.

\vspace{0.5cm}
\noindent\rule{\textwidth}{0.4pt}
\vspace{0.5cm}

\noindent
\textbf{Contact Information:} \\
Michael Kernaghan \\
Director, Abbotsford Croquet Club \\
michael.kernaghan@gmail.com \\
604-615-9138

\end{document}
